% GENERATED FILE. Edit canonical lessons, then run npm run generate:books.
% canonical-source-hash: fnv1a64:ecf7b83211d42d01
% canonical-lessons: ZH-C80-shangban, ZH-C80-gongren, ZH-C80-siji, ZH-C80-yisheng, ZH-C80-gongsi

\chapter{Jobs}
\label{ch:zh-jobs}
\hlchaptermodality{80}

\begin{chapteropening}
\textbf{By the end of this chapter:} \emph{I can say who goes to work where: workers, drivers, doctors and companies.}

Complete the last lesson of chapter 80: I can say who goes to work where: workers, drivers, doctors and companies.
\end{chapteropening}

\section[shàngbān]{shàngbān — to go to work}

\label{lesson:ZH-C80-shangban}



\begin{quote}
Before the new one: say the Mandarin for homework, then the Mandarin for to take an exam.
\end{quote}

\subsection*{You'll want to know: shàngbān}
\emph{shàngbān} — \textquotedblleft{}to go to work\textquotedblright{}.

\begin{quote}
\emph{wǒ jiǔ diǎn shàngbān} — I start work at nine.
\end{quote}

Its characters come later; for now it is a word for the ear.

\begin{cousinweb}
\emph{shàngbān} is built from \zh{上} 'to go up to' and bān, both inherited from Middle Chinese, \emph{dzyangX} and \emph{paen} (Guangyun, in Baxter's transcription). The dictionary explains the second, through an older book, as dividing up jade tallies; a shift of work is a likely step, but it is not settled.
\end{cousinweb}

\begin{grammarlens}[title={What you've built}]
One more word about people.
\end{grammarlens}

\subsection*{Your turn}
\begin{itemize}
\raggedright
  \item \emph{Say it:} \emph{shàngbān}
  \item \emph{Say it:} \emph{shàngbān}, once more
  \item \emph{Say it:} \emph{wǒ jiǔ diǎn shàngbān} — I start work at nine.
  \item \emph{Recall:} say the Mandarin for homework, then the Mandarin for to take an exam, then say \emph{shàngbān} again
\end{itemize}

\begin{tcolorbox}[breakable,colback=teal!4,colframe=teal!35!black,title={Before you move on}]
What does shàngbān mean? (\textquotedblleft{}To go to work\textquotedblright{}.) Say it once more.
\end{tcolorbox}
\section[gōngrén]{gōngrén — worker}

\label{lesson:ZH-C80-gongren}



\begin{quote}
Before the new one: say the Mandarin for to take an exam, then the Mandarin for to go to work.
\end{quote}

\subsection*{You'll want to know: gōngrén}
\emph{gōngrén} — \textquotedblleft{}worker\textquotedblright{}.

\begin{quote}
\emph{wǒ bàba shì gōngrén} — My dad is a worker.
\end{quote}

Its characters come later; for now it is a word for the ear.

\begin{cousinweb}
\emph{gōngrén} is built from gōng 'craft, work' and \zh{人} 'person', both inherited from Middle Chinese, \emph{kuwng} and \emph{nyin} (Guangyun, in Baxter's transcription).
\end{cousinweb}

\begin{grammarlens}[title={What you've built}]
One more word about people.
\end{grammarlens}

\subsection*{Your turn}
\begin{itemize}
\raggedright
  \item \emph{Say it:} \emph{gōngrén}
  \item \emph{Say it:} \emph{gōngrén}, once more
  \item \emph{Say it:} \emph{wǒ bàba shì gōngrén} — My dad is a worker.
  \item \emph{Recall:} say the Mandarin for to take an exam, then the Mandarin for to go to work, then say \emph{gōngrén} again
\end{itemize}

\begin{tcolorbox}[breakable,colback=teal!4,colframe=teal!35!black,title={Before you move on}]
What does gōngrén mean? (\textquotedblleft{}Worker\textquotedblright{}.) Say it once more.
\end{tcolorbox}
\section[sījī]{sījī — driver}

\label{lesson:ZH-C80-siji}



\begin{quote}
Before the new one: say the Mandarin for to go to work, then the Mandarin for worker.
\end{quote}

\subsection*{You'll want to know: sījī}
\emph{sījī} — \textquotedblleft{}driver\textquotedblright{}.

\begin{quote}
\emph{tā shì chūzūchē sījī} — He's a taxi driver.
\end{quote}

Its characters come later; for now it is a word for the ear.

\begin{cousinweb}
\emph{sījī} is built from sī 'to be in charge' and \zh{机} 'machine', both inherited from Middle Chinese, \emph{si} and \emph{kj+j} (Guangyun, in Baxter's transcription): the one in charge of the machine.
\end{cousinweb}

\begin{grammarlens}[title={What you've built}]
One more word about people.
\end{grammarlens}

\subsection*{Your turn}
\begin{itemize}
\raggedright
  \item \emph{Say it:} \emph{sījī}
  \item \emph{Say it:} \emph{sījī}, once more
  \item \emph{Say it:} \emph{tā shì chūzūchē sījī} — He's a taxi driver.
  \item \emph{Recall:} say the Mandarin for to go to work, then the Mandarin for worker, then say \emph{sījī} again
\end{itemize}

\begin{tcolorbox}[breakable,colback=teal!4,colframe=teal!35!black,title={Before you move on}]
What does sījī mean? (\textquotedblleft{}Driver\textquotedblright{}.) Say it once more.
\end{tcolorbox}
\section[yīshēng]{yīshēng — doctor}

\label{lesson:ZH-C80-yisheng}



\begin{quote}
Before the new one: say the Mandarin for worker, then the Mandarin for driver.
\end{quote}

\subsection*{You'll want to know: yīshēng}
\emph{yīshēng} — \textquotedblleft{}doctor\textquotedblright{}.

\begin{quote}
\emph{wǒ māma shì yīshēng} — My mum is a doctor.
\end{quote}

Its characters come later; for now it is a word for the ear.

\begin{cousinweb}
\emph{yīshēng} is built from yī 'healing' and \zh{生}, as in \zh{学生} and in xiānsheng, both inherited from Middle Chinese, \emph{'i} and \emph{srjaeng} (Guangyun, in Baxter's transcription).
\end{cousinweb}

\begin{grammarlens}[title={What you've built}]
One more word about people.
\end{grammarlens}

\subsection*{Your turn}
\begin{itemize}
\raggedright
  \item \emph{Say it:} \emph{yīshēng}
  \item \emph{Say it:} \emph{yīshēng}, once more
  \item \emph{Say it:} \emph{wǒ māma shì yīshēng} — My mum is a doctor.
  \item \emph{Recall:} say the Mandarin for worker, then the Mandarin for driver, then say \emph{yīshēng} again
\end{itemize}

\begin{tcolorbox}[breakable,colback=teal!4,colframe=teal!35!black,title={Before you move on}]
What does yīshēng mean? (\textquotedblleft{}Doctor\textquotedblright{}.) Say it once more.
\end{tcolorbox}
\section[gōngsī]{gōngsī — company; firm}

\label{lesson:ZH-C80-gongsi}



\begin{quote}
Before the new one: say the Mandarin for driver, then the Mandarin for doctor.
\end{quote}

\subsection*{You'll want to know: gōngsī}
\emph{gōngsī} — \textquotedblleft{}company; firm\textquotedblright{}.

\begin{quote}
\emph{wǒ de gōngsī hěn dà} — My company is big.
\end{quote}

Its characters come later; for now it is a word for the ear.

\begin{cousinweb}
\emph{gōngsī} is built from gōng 'public, shared' and sī 'to manage', both inherited from Middle Chinese, \emph{kuwng} and \emph{si} (Guangyun, in Baxter's transcription).
\end{cousinweb}

\begin{grammarlens}[title={What you've built}]
One more word about people.
\end{grammarlens}

\subsection*{Your turn}
\begin{itemize}
\raggedright
  \item \emph{Say it:} \emph{gōngsī}
  \item \emph{Say it:} \emph{gōngsī}, once more
  \item \emph{Say it:} \emph{wǒ de gōngsī hěn dà} — My company is big.
  \item \emph{Recall:} say the Mandarin for driver, then the Mandarin for doctor, then say \emph{gōngsī} again
\end{itemize}

\begin{tcolorbox}[breakable,colback=teal!4,colframe=teal!35!black,title={Before you move on}]
What does gōngsī mean? (\textquotedblleft{}Company; firm\textquotedblright{}.) Say it once more.
\end{tcolorbox}
